\documentclass[11pt]{article}
\usepackage[margin=1in]{geometry}
\usepackage{amsmath,amssymb,amsthm,mathtools}
\usepackage{graphicx}
\usepackage{microtype}
\usepackage[hidelinks]{hyperref}
\usepackage{enumitem}

\newtheorem{theorem}{Theorem}
\newtheorem{corollary}{Corollary}
\newcommand{\T}{\mathbb T}
\newcommand{\N}{\mathbb N}
\newcommand{\E}{\mathcal E}
\newcommand{\rank}{\operatorname{rank}}
\newcommand{\diag}{\operatorname{diag}}

\title{Cowen-Set Membership Does Not Determine\\
the Blaschke--Potapov Degree}
\author{Automated mathematics run \texttt{fa\_banach\_001}, lane 7}
\date{August 11, 2026}

\begin{document}
\maketitle

\begin{center}
\fbox{\parbox{0.9\textwidth}{\textbf{Status.}
Candidate full counterexample, likely valid, pending expert review.
Problem~4.2 of Abhinand--Curto--Prasad has a negative answer, already for
analytic diagonal monomial symbols.  For this family we also determine the
exact set of determinant degrees occurring in the matrix Cowen set.}}
\end{center}

\section{Source question and answer}

For $\Phi\in L^\infty(\T,M_n)$, the source defines
\[
 \E(\Phi)=\{K\in H^\infty(\T,M_n):\|K\|_\infty\leq1
       \text{ and }\Phi-K\Phi^*\in H^\infty(\T,M_n)\}.
\]
Problem~4.2 asks whether, for a normal symbol $\Phi$, the presence of a finite
Blaschke--Potapov product $B\in\E(\Phi)$ forces
\[
 \rank[T_\Phi^*,T_\Phi]=\deg(\det B).
\]

\begin{figure}[ht]
  \centering
  \includegraphics[width=0.98\textwidth]{figures/open_problem_crop.png}
  \caption{Problem~4.2 on page 11 of arXiv:2605.02214.}
  \label{fig:problem}
\end{figure}

The answer is no.  Indeed, for any $m,n\geq1$ and $k>m$, take
\[
 \Phi_m(z)=z^m I_n,\qquad B_k(z)=z^k I_n.
\]
Then $B_k\in\E(\Phi_m)$, but the two quantities in the proposed identity are
$mn$ and $kn$.  Thus one may take the concrete block example
$\Phi(z)=zI_2$, $B(z)=z^2I_2$, giving ranks/degrees $2$ and $4$.

\subsection*{Proof intuition}

The Cowen-set condition is stable under adding extra inner factors on the
analytic side.  For a monomial symbol $z^mI_n$, every $z^kI_n$ with $k\geq m$
still makes $\Phi-B\Phi^*$ analytic.  The Toeplitz commutator, however, only
sees the fixed symbol and has rank $mn$.  The proposed equality therefore
cannot hold for an arbitrary member $B$ of the Cowen set.  More sharply, for
diagonal monomials, analyticity forces a minimal inner divisor; all additional
inner degree is optional.

\section{A full counterexample}

Let $S$ denote the unilateral shift on $H^2(\T)$.  We use the convention
$[A^*,A]=A^*A-AA^*$.

\begin{theorem}[Monomial counterexample]\label{thm:counterexample}
For every $m,n\geq1$ and every $k>m$, the pointwise normal symbol
$\Phi_m(z)=z^mI_n$ and the finite Blaschke--Potapov product $B_k(z)=z^kI_n$
satisfy
\[
 B_k\in\E(\Phi_m),\qquad
 \rank[T_{\Phi_m}^*,T_{\Phi_m}]=mn,\qquad
 \deg(\det B_k)=kn.
\]
Consequently Problem~4.2 is false, including when $n\geq2$.
\end{theorem}

\begin{proof}
On $\T$ we have $\Phi_m^*(z)=z^{-m}I_n$, and hence
\[
 \Phi_m-B_k\Phi_m^*=(z^m-z^{k-m})I_n\in H^\infty(\T,M_n).
\]
The matrix $B_k$ is inner, so $\|B_k\|_\infty=1$; this proves
$B_k\in\E(\Phi_m)$.  Pointwise, $\Phi_m$ is a scalar multiple of the identity
and is therefore normal (in fact unitary on $\T$).

Under $H^2(\T,\mathbb C^n)\cong H^2(\T)\otimes\mathbb C^n$,
\[
 T_{\Phi_m}=S^m\otimes I_n.
\]
Since $S^{*m}S^m=I$ and $I-S^mS^{*m}$ is the projection onto
$\operatorname{span}\{1,z,\ldots,z^{m-1}\}$,
\[
 [T_{\Phi_m}^*,T_{\Phi_m}]
   =(I-S^mS^{*m})\otimes I_n
\]
has rank $mn$.  Finally $\det B_k=z^{kn}$, whose degree is $kn$.  Because
$k>m$, the proposed equality fails.
\end{proof}

\section{Exact degree spectrum for diagonal monomials}

The same mechanism yields a sharp result, not merely one counterexample.

\begin{theorem}[Minimum-degree theorem]\label{thm:spectrum}
Let
\[
 D(z)=\diag(z^{m_1},\ldots,z^{m_n}),\qquad m_j\geq1,
 \qquad r=\sum_{j=1}^n m_j.
\]
Then
\[
 \rank[T_D^*,T_D]=r
\]
and the finite Blaschke--Potapov products in $\E(D)$ have exactly the
determinant degrees
\[
 \bigl\{\deg(\det B):B\text{ finite Blaschke--Potapov},\ B\in\E(D)\bigr\}
   =\{r,r+1,r+2,\ldots\}.
\]
In particular, the commutator rank is the \emph{minimum}, not the common
value, of the determinant degrees in the Cowen set.
\end{theorem}

\begin{proof}
The Toeplitz operator is the orthogonal direct sum
$T_D=\bigoplus_{j=1}^n S^{m_j}$.  The computation in
Theorem~\ref{thm:counterexample} gives
\[
 [T_D^*,T_D]=\bigoplus_{j=1}^n (I-S^{m_j}S^{*m_j}),
\]
so its rank is $r$.

Now let a finite Blaschke--Potapov product $B=(B_{ij})$ belong to $\E(D)$.
The $(i,j)$ entry of $BD^*$ on $\T$ is $B_{ij}(z)z^{-m_j}$.  Since
$D-BD^*$ has no negative Fourier coefficients, neither can this entry: the
analytic entry $D_{ij}$ has none to cancel.  Thus every Taylor coefficient of
$B_{ij}$ below order $m_j$ vanishes.  Equivalently, the $j$th column of $B$
is divisible by $z^{m_j}$, and
\[
 B=CD
\]
for an analytic rational matrix function $C$.

Both $B$ and $D$ are unitary almost everywhere on $\T$, so
$C=BD^*$ is unitary there.  Hence $C$ is rational inner.  Taking
determinants gives
\[
 \det B=(\det C)z^r.
\]
The scalar rational inner function $\det C$ has a nonnegative integer degree;
therefore every such $B$ satisfies $\deg(\det B)\geq r$.

Conversely, for each $d\geq0$ set
\[
 B_d(z)=\diag(z^{m_1+d},z^{m_2},\ldots,z^{m_n}).
\]
This is a finite Blaschke--Potapov product and
\[
 B_dD^*=\diag(z^d,1,\ldots,1),
\]
so $D-B_dD^*$ is analytic and $B_d\in\E(D)$.  Moreover
$\deg(\det B_d)=r+d$.  This realizes every integer at least $r$.
\end{proof}

\section{Quantifier audit and sharp repair}

Problem~4.2 names a particular product $B$ already known to lie in
$\E(\Phi)$ and asks whether its degree must equal the rank.  Theorems
\ref{thm:counterexample} and~\ref{thm:spectrum} refute exactly that literal
statement.  There is no ambiguity from block size: the examples work for all
$n$, and in particular for every $n\geq2$.

The paper's central conjecture has a different, existential quantifier: under
hyponormality and finite-rank hypotheses, it asks whether there \emph{exists}
a finite Blaschke--Potapov product in $\E(\Phi)$ whose determinant degree is
the commutator rank.  Our examples do not refute that conjecture.  For the
diagonal symbols above, $B=D$ belongs to $\E(D)$ and has degree $r$, so it is
exactly the desired witness.

Thus the sharp repaired statement for this natural class is
\[
 \rank[T_D^*,T_D]
 =\min\{\deg(\det B):B\text{ finite Blaschke--Potapov in }\E(D)\},
\]
whereas equality for every available $B$ is false.  Whether this
minimum-degree formulation holds for broader normal symbols is not claimed.

\section{Verification, scope, and novelty}

The proof uses only Fourier coefficients and the standard Hardy basis.  The
review-critical checks are: the sign of the self-commutator, the right-column
exponents in $BD^*$, and the universal-versus-existential quantifier.  All
three computations are displayed explicitly.

This packet completely answers Problem~4.2 negatively and gives an exact
classification for diagonal monomial symbols.  It does not settle the
paper's central existential conjecture for arbitrary symbols.

Local run indexes were searched by arXiv identifier, title, and the core
Cowen-set/degree terminology.  Bounded primary-source searches through
August~11, 2026 found the source paper but no later answer or correction.
The argument is elementary and may be folklore; novelty is asserted only as
an explicit answer to this 2026 problem and remains subject to expert review.

\begin{thebibliography}{9}
\bibitem{source}
M.~Abhinand, R.~E.~Curto, and T.~Prasad,
\emph{Hyponormal block Toeplitz operators with finite rank self-commutators},
arXiv:2605.02214 (2026), Problem~4.2, p.~11.
\end{thebibliography}

\end{document}
